\chapter{Complexity — CNF}
%%% generated by the blueprint pipeline from TCSlib.Complexity.Formulas.CNF
%%%%%%%%%%%%%%%%%%%%%%%%%%%%%%%%%%%%%%%%%%%%%%%%%%%%%%%%%%%%%%%%%%%%%%%%%%%%%%%
\section{Overview}

This module supplies the conjunctive-normal-form layer used by the development of
$\mathrm{SAT}$, $3\mathrm{SAT}$ and the Cook--Levin theorem [AB09, Section~2.3.1]:
satisfiability of a CNF formula, the variable count of a formula, clause-width bounds
($k$CNF), the fact that evaluation reads only the variables below the variable count,
and CNF universality [AB09, Claim~2.13], namely that every Boolean function on $\ell$
bits is computed by a CNF formula with at most $2^\ell$ clauses, each of width at
most $\ell$.

Throughout this chapter a CNF formula is a finite list of clauses, each clause a
finite list of literals; a literal $x_v^b$ consists of a variable $v$ and a polarity
$b \in \{0,1\}$ and is satisfied by an assignment $a$ exactly when $a(v) = b$. A
formula evaluates to true under $a$ when every clause contains a satisfied literal, so
the formula with no clauses is true and an empty clause is false.

%%%%%%%%%%%%%%%%%%%%%%%%%%%%%%%%%%%%%%%%%%%%%%%%%%%%%%%%%%%%%%%%%%%%%%%%%%%%%%%
\section{Declarations}

\begin{definition}[Satisfiable CNF formula]
\lean{Std.Sat.CNF.Satisfiable}
\label{Std.Sat.CNF.Satisfiable}
\leanok
Let $\alpha$ be a type of variables and let $\varphi$ be a CNF formula over $\alpha$,
that is, a conjunction of clauses, each a disjunction of literals $x_v^b$
($v \in \alpha$, $b \in \{0,1\}$), where $x_v^b$ is satisfied by an assignment $a$
exactly when $a(v) = b$. The formula $\varphi$ is \emph{satisfiable} if there exists
an assignment $a : \alpha \to \{0,1\}$ under which $\varphi$ evaluates to true
[AB09, Section~2.3.1].
\end{definition}

\begin{definition}[Variable count of a CNF formula]
\lean{Std.Sat.CNF.numVars}
\label{Std.Sat.CNF.numVars}
\leanok
Let $\varphi$ be a CNF formula whose variables are indexed by natural numbers. Its
\emph{variable count} $N(\varphi)$ is one plus the largest index of a variable
occurring in some literal of $\varphi$, and $0$ if no literal occurs in $\varphi$ (for
instance when $\varphi$ has no clauses, or only empty clauses). Equivalently,
$N(\varphi)$ is the least $N \in \mathbb{N}$ such that every variable occurring in
$\varphi$ has index less than $N$.
\end{definition}

\begin{definition}[Clause width at most $k$ ($k$CNF)]
\lean{Std.Sat.CNF.WidthAtMost}
\label{Std.Sat.CNF.WidthAtMost}
\leanok
Let $\alpha$ be a type of variables, let $\varphi$ be a CNF formula over $\alpha$, and
let $k \in \mathbb{N}$. The formula $\varphi$ has \emph{width at most $k$} (is a
$k$CNF) if every clause of $\varphi$ contains at most $k$ literals
[AB09, Section~2.3.1]. In particular, a formula with no clauses has width at most $k$
for every $k$.
\end{definition}

\begin{theorem}[Members are bounded by the list maximum]
\lean{Complexity.le_foldr_max_of_mem}
\label{Complexity.le_foldr_max_of_mem}
\leanok
\difficulty{4}
For a finite list $s$ of natural numbers, write $\max(s)$ for the right fold of
$\max$ over $s$ with initial value $0$ (so the empty list has maximum $0$). If $n$ is a
member of $s$, then $n \le \max(s)$.
\end{theorem}

\begin{theorem}[Evaluation reads only variables below the count]
\lean{Complexity.eval_congr_of_lt_numVars}
\label{Complexity.eval_congr_of_lt_numVars}
\leanok
\uses{Complexity.le_foldr_max_of_mem, Std.Sat.CNF.numVars}
\difficulty{4}
Let $\varphi$ be a CNF formula whose variables are indexed by natural numbers, and let
$N(\varphi)$ be one plus the largest variable index occurring in $\varphi$ ($0$ if no
variable occurs). If two assignments $a, b : \mathbb{N} \to \{0,1\}$ satisfy
$a(v) = b(v)$ for every $v < N(\varphi)$, then $\varphi$ evaluates to the same truth
value under $a$ as under $b$.
\end{theorem}

\begin{theorem}[Termwise bound bounds the list maximum]
\lean{Complexity.foldr_max_le_of_forall}
\label{Complexity.foldr_max_le_of_forall}
\leanok
\uses{Complexity.falsifyingClause_eval_false}
\difficulty{4}
Let $s$ be a finite list of natural numbers, write $\max(s)$ for the right fold of
$\max$ over $s$ with initial value $0$, and let $n \in \mathbb{N}$. If every member of
$s$ is at most $n$, then $\max(s) \le n$; in particular this holds when $s$ is empty.
\end{theorem}

\begin{definition}[Clause excluding a single assignment]
\lean{Complexity.falsifyingClause}
\label{Complexity.falsifyingClause}
\leanok
Let $\ell \in \mathbb{N}$ and $v = (v_0, \dots, v_{\ell-1}) \in \{0,1\}^\ell$. The
clause \emph{excluding} $v$ is
\[ C_v = \bigvee_{i=0}^{\ell-1} x_i^{\neg v_i}, \]
the disjunction, in order of $i$, of $\ell$ literals on the variables
$x_0, \dots, x_{\ell-1}$, where the literal $x_i^{\neg v_i}$ is satisfied by an
assignment $a : \mathbb{N} \to \{0,1\}$ exactly when $a(i) \ne v_i$
[AB09, Claim~2.13].
\end{definition}

\begin{theorem}[Excluding clause is false exactly at its point]
\lean{Complexity.falsifyingClause_eval_false}
\label{Complexity.falsifyingClause_eval_false}
\leanok
\uses{Complexity.falsifyingClause}
\difficulty{4}
Let $\ell \in \mathbb{N}$, let $v \in \{0,1\}^\ell$, and let
$C_v = \bigvee_{i<\ell} x_i^{\neg v_i}$ be the clause on the variables
$x_0, \dots, x_{\ell-1}$ whose $i$-th literal is satisfied by an assignment $a$ exactly
when $a(i) \ne v_i$. For every assignment $a : \mathbb{N} \to \{0,1\}$, the clause
$C_v$ evaluates to false under $a$ if and only if $a(i) = v_i$ for all $i < \ell$
[AB09, Claim~2.13].
\end{theorem}

\begin{definition}[Truth-table CNF of a Boolean function]
\lean{Complexity.falsifyingCNF}
\label{Complexity.falsifyingCNF}
\leanok
\uses{Complexity.falsifyingClause}
Let $\ell \in \mathbb{N}$ and let $f : \{0,1\}^\ell \to \{0,1\}$. The
\emph{truth-table CNF} $\varphi_f$ of $f$ is the conjunction, over all
$v \in \{0,1\}^\ell$ with $f(v) = 0$, of the clause
$C_v = \bigvee_{i<\ell} x_i^{\neg v_i}$, whose $i$-th literal is satisfied by an
assignment $a : \mathbb{N} \to \{0,1\}$ exactly when $a(i) \ne v_i$
[AB09, Claim~2.13].
\end{definition}

\begin{theorem}[Truth-table CNF uses only the first $\ell$ variables]
\lean{Complexity.falsifyingCNF_numVars}
\label{Complexity.falsifyingCNF_numVars}
\leanok
\uses{Complexity.falsifyingCNF, Complexity.foldr_max_le_of_forall, Std.Sat.CNF.numVars}
\difficulty{3}
Let $\ell \in \mathbb{N}$, let $f : \{0,1\}^\ell \to \{0,1\}$, and let $\varphi_f$ be
the CNF formula that is the conjunction, over all $v \in \{0,1\}^\ell$ with
$f(v) = 0$, of the clause $\bigvee_{i<\ell} x_i^{\neg v_i}$ whose $i$-th literal is
satisfied by an assignment $a$ exactly when $a(i) \ne v_i$. Then every variable
occurring in $\varphi_f$ has index less than $\ell$; that is, one plus the largest
variable index occurring in $\varphi_f$ (or $0$ if none occurs) is at most $\ell$.
\end{theorem}

\begin{theorem}[Truth-table CNF has at most $2^\ell$ clauses]
\lean{Complexity.falsifyingCNF_length}
\label{Complexity.falsifyingCNF_length}
\leanok
\uses{Complexity.falsifyingCNF}
\difficulty{2}
Let $\ell \in \mathbb{N}$, let $f : \{0,1\}^\ell \to \{0,1\}$, and let $\varphi_f$ be
the CNF formula that is the conjunction, over all $v \in \{0,1\}^\ell$ with
$f(v) = 0$, of the clause $\bigvee_{i<\ell} x_i^{\neg v_i}$ whose $i$-th literal is
satisfied by an assignment $a$ exactly when $a(i) \ne v_i$. Then $\varphi_f$ has at
most $2^\ell$ clauses.
\end{theorem}

\begin{theorem}[Truth-table CNF has width at most $\ell$]
\lean{Complexity.falsifyingCNF_width}
\label{Complexity.falsifyingCNF_width}
\leanok
\uses{Complexity.falsifyingCNF, Std.Sat.CNF.WidthAtMost}
\difficulty{3}
Let $\ell \in \mathbb{N}$, let $f : \{0,1\}^\ell \to \{0,1\}$, and let $\varphi_f$ be
the CNF formula that is the conjunction, over all $v \in \{0,1\}^\ell$ with
$f(v) = 0$, of the clause $\bigvee_{i<\ell} x_i^{\neg v_i}$ whose $i$-th literal is
satisfied by an assignment $a$ exactly when $a(i) \ne v_i$. Then every clause of
$\varphi_f$ contains at most $\ell$ literals, i.e.\ $\varphi_f$ is an $\ell$CNF.
\end{theorem}

\begin{theorem}[Truth-table CNF computes its function]
\lean{Complexity.falsifyingCNF_eval}
\label{Complexity.falsifyingCNF_eval}
\leanok
\uses{Complexity.falsifyingCNF, Complexity.falsifyingCNF_width, Complexity.falsifyingClause, Complexity.falsifyingClause_eval_false}
\difficulty{5}
Let $\ell \in \mathbb{N}$, let $f : \{0,1\}^\ell \to \{0,1\}$, and let $\varphi_f$ be
the CNF formula that is the conjunction, over all $v \in \{0,1\}^\ell$ with
$f(v) = 0$, of the clause $\bigvee_{i<\ell} x_i^{\neg v_i}$ whose $i$-th literal is
satisfied by an assignment $a$ exactly when $a(i) \ne v_i$. Then for every assignment
$a : \mathbb{N} \to \{0,1\}$, the formula $\varphi_f$ evaluates under $a$ to
$f(a(0), a(1), \dots, a(\ell-1))$.
\end{theorem}

\begin{theorem}[CNF universality]
\lean{Complexity.exists_cnf_boolFun}
\label{Complexity.exists_cnf_boolFun}
\leanok
\uses{Complexity.falsifyingCNF, Complexity.falsifyingCNF_eval, Complexity.falsifyingCNF_length, Complexity.falsifyingCNF_numVars, Complexity.falsifyingCNF_width, Std.Sat.CNF.WidthAtMost, Std.Sat.CNF.numVars}
\difficulty{1}
For every $\ell \in \mathbb{N}$ and every Boolean function
$f : \{0,1\}^\ell \to \{0,1\}$ there exists a CNF formula $\varphi$ over variables
indexed by natural numbers such that: every variable occurring in $\varphi$ has index
less than $\ell$ (one plus the largest occurring index, or $0$ if none occurs, is at
most $\ell$); $\varphi$ has at most $2^\ell$ clauses; every clause of $\varphi$
contains at most $\ell$ literals; and for every assignment
$a : \mathbb{N} \to \{0,1\}$, the formula $\varphi$ evaluates under $a$ to
$f(a(0), a(1), \dots, a(\ell-1))$ [AB09, Claim~2.13].
\end{theorem}
